%%%%%%%%%%%%%%%%%%%%%%%%%%%%%%%%%%%%%%%%%%%%%%%%%%%%%%%%%%%%%%%
%% OXFORD THESIS TEMPLATE

% Use this template to produce a standard thesis that meets the Oxford University requirements for DPhil submission
%
% Originally by Keith A. Gillow (gillow@maths.ox.ac.uk), 1997
% Modified by Sam Evans (sam@samuelevansresearch.org), 2007
% Modified by John McManigle (john@oxfordechoes.com), 2015
%
% This version Copyright (c) 2015-2023 John McManigle
%
% Broad permissions are granted to use, modify, and distribute this software
% as specified in the MIT License included in this distribution's LICENSE file.
%

% I've (John) tried to comment this file extensively, so read through it to see how to use the various options.  Remember
% that in LaTeX, any line starting with a % is NOT executed.  Several places below, you have a choice of which line to use
% out of multiple options (eg draft vs final, for PDF vs for binding, etc.)  When you pick one, add a % to the beginning of
% the lines you don't want.


%%%%% CHOOSE PAGE LAYOUT
% The most common choices should be below.  You can also do other things, like replacing "a4paper" with "letterpaper", etc.

% This one will format for two-sided binding (ie left and right pages have mirror margins; blank pages inserted where needed):
\documentclass[a4paper,twoside]{ociamthesis}
% This one will format for one-sided binding (ie left margin > right margin; no extra blank pages):
%\documentclass[a4paper]{ociamthesis}
% This one will format for PDF output (ie equal margins, no extra blank pages):
%\documentclass[a4paper,nobind]{ociamthesis} 



%%%%% SELECT YOUR DRAFT OPTIONS
% Three options going on here; use in any combination.  But remember to turn the first two off before
% generating a PDF to send to the printer!

% This adds a "DRAFT" footer to every normal page.  (The first page of each chapter is not a "normal" page.)
\fancyfoot[C]{\emph{DRAFT Printed on \today}}  

% This highlights (in blue) corrections marked with (for words) \mccorrect{blah} or (for whole
% paragraphs) \begin{mccorrection} . . . \end{mccorrection}.  This can be useful for sending a PDF of
% your corrected thesis to your examiners for review.  Turn it off, and the blue disappears.
\correctionstrue


%%%%% BIBLIOGRAPHY SETUP
% Note that your bibliography will require some tweaking depending on your department, preferred format, etc.
% The options included below are just very basic "sciencey" and "humanitiesey" options to get started.
% If you've not used LaTeX before, I recommend reading a little about biblatex/biber and getting started with it.
% If you're already a LaTeX pro and are used to natbib or something, modify as necessary.
% Either way, you'll have to choose and configure an appropriate bibliography format...

% The science-type option: numerical in-text citation with references in order of appearance.
\usepackage[style=numeric-comp, sorting=none, backend=biber, doi=false, isbn=false]{biblatex}
\newcommand*{\bibtitle}{References}

% The humanities-type option: author-year in-text citation with an alphabetical works cited.
%\usepackage[style=authoryear, sorting=nyt, backend=biber, maxcitenames=2, useprefix, doi=false, isbn=false]{biblatex}
%\newcommand*{\bibtitle}{Works Cited}

% This makes the bibliography left-aligned (not 'justified') and slightly smaller font.
\renewcommand*{\bibfont}{\raggedright\small}

% Change this to the name of your .bib file (usually exported from a citation manager like Zotero or EndNote).
\addbibresource{references.bib}


% Uncomment this if you want equation numbers per section (2.3.12), instead of per chapter (2.18):
%\numberwithin{equation}{subsection}



%%%%% THESIS / TITLE PAGE INFORMATION
% Everybody needs to complete the following:
\title{Scaling Up Quantum Compilation}
\author{Luca M Mondada}
\college{Lady Margaret Hall}

% Master's candidates who require the alternate title page (with candidate number and word count)
% must also un-comment and complete the following three lines:
%\masterssubmissiontrue
%\candidateno{933516}
%\wordcount{28,815}

% Uncomment the following line if your degree also includes exams (eg most masters):
%\renewcommand{\submittedtext}{Submitted in partial completion of the}
% Your full degree name.  (But remember that DPhils aren't "in" anything.  They're just DPhils.)
\degree{Doctor of Philosophy}
% Term and year of submission, or date if your board requires (eg most masters)
\degreedate{Michaelmas 2024}


%%%%% YOUR OWN PERSONAL MACROS / stuff
% This is a good place to dump your own LaTeX macros as they come up.
\usepackage{hyperref}
\usepackage{cleveref}

% To make text superscripts shortcuts
\renewcommand{\th}{\textsuperscript{th}} % ex: I won 4\th place
\newcommand{\nd}{\textsuperscript{nd}}
\renewcommand{\st}{\textsuperscript{st}}
\newcommand{\rd}{\textsuperscript{rd}}

%%%%% THE ACTUAL DOCUMENT STARTS HERE
\begin{document}



%%%%% CHOOSE YOUR LINE SPACING HERE
% This is the official option.  Use it for your submission copy and library copy:
% \setlength{\textbaselineskip}{22pt plus2pt}
% This is closer spacing (about 1.5-spaced) that you might prefer for your personal copies:
\setlength{\textbaselineskip}{18pt plus2pt minus1pt}

% You can set the spacing here for the roman-numbered pages (acknowledgements, table of contents, etc.)
\setlength{\frontmatterbaselineskip}{17pt plus1pt minus1pt}

% Leave this line alone; it gets things started for the real document.
\setlength{\baselineskip}{\textbaselineskip}


%%%%% CHOOSE YOUR SECTION NUMBERING DEPTH HERE
% You have two choices.  First, how far down are sections numbered?  (Below that, they're named but
% don't get numbers.)  Second, what level of section appears in the table of contents?  These don't have
% to match: you can have numbered sections that don't show up in the ToC, or unnumbered sections that
% do.  Throughout, 0 = chapter; 1 = section; 2 = subsection; 3 = subsubsection, 4 = paragraph...

% The level that gets a number:
\setcounter{secnumdepth}{2}
% The level that shows up in the ToC:
\setcounter{tocdepth}{2}


%%%%% ABSTRACT SEPARATE
% This is used to create the separate, one-page abstract that you are required to hand into the Exam
% Schools.  You can comment it out to generate a PDF for printing or whatnot.
\begin{abstractseparate}
	Your abstract text goes here.  Check your departmental regulations, but generally this should be less than 300 words.  See the beginning of Chapter~\ref{ch:2-litreview} for more.

Lorem ipsum dolor sit amet, consectetur adipiscing elit. Pellentesque sit amet nibh volutpat, scelerisque nibh a, vehicula neque. Integer placerat nulla massa, et vestibulum velit dignissim id. Ut eget nisi elementum, consectetur nibh in, condimentum velit. Quisque sodales dui ut tempus mattis. Duis malesuada arcu at ligula egestas egestas. Phasellus interdum odio at sapien fringilla scelerisque. Mauris sagittis eleifend sapien, sit amet laoreet felis mollis quis. Pellentesque dui ante, finibus eget blandit sit amet, tincidunt eu neque. Vivamus rutrum dapibus ligula, ut imperdiet lectus tincidunt ac. Pellentesque ac lorem sed diam egestas lobortis.

Suspendisse leo purus, efficitur mattis urna a, maximus molestie nisl. Aenean porta semper tortor a vestibulum. Suspendisse viverra facilisis lorem, non pretium erat lacinia a. Vestibulum tempus, quam vitae placerat porta, magna risus euismod purus, in viverra lorem dui at metus. Sed ac sollicitudin nunc. In maximus ipsum nunc, placerat maximus tortor gravida varius. Suspendisse pretium, lorem at porttitor rhoncus, nulla urna condimentum tortor, sed suscipit nisi metus ac risus.

Aenean sit amet enim quis lorem tristique commodo vitae ut lorem. Duis vel tincidunt lacus. Sed massa velit, lacinia sed posuere vitae, malesuada vel ante. Praesent a rhoncus leo. Etiam sed rutrum enim. Pellentesque lobortis elementum augue, at suscipit justo malesuada at. Lorem ipsum dolor sit amet, consectetur adipiscing elit. Praesent rhoncus convallis ex. Etiam commodo nunc ex, non consequat diam consectetur ut. Pellentesque vitae est nec enim interdum dapibus. Donec dapibus purus ipsum, eget tincidunt ex gravida eget. Donec luctus nisi eu fringilla mollis. Donec eget lobortis diam.

Suspendisse finibus placerat dolor. Etiam ornare elementum ex ut vehicula. Donec accumsan mattis erat. Quisque cursus fringilla diam, eget placerat neque bibendum eu. Ut faucibus dui vitae dolor porta, at elementum ipsum semper. Sed ultrices dui non arcu pellentesque placerat. Etiam posuere malesuada turpis, nec malesuada tellus malesuada. % Create an abstract.tex file in the 'text' folder for your abstract.
\end{abstractseparate}


% JEM: Pages are roman numbered from here, though page numbers are invisible until ToC.  This is in
% keeping with most typesetting conventions.
\begin{romanpages}

% JEM: By default, this template uses the traditional Oxford "Belt Crest". Un-comment the following
% line to use the newer, "Blue Square" logo:
% \renewcommand{\crest}{{\includegraphics[width=4.2cm, height=4.2cm]{figures/newlogo.pdf}}}

% Title page is created here
\maketitle

%%%%% DEDICATION -- If you'd like one, un-comment the following.
%\begin{dedication}
%This thesis is dedicated to\\
%someone\\
%for some special reason\\
%\end{dedication}

%%%%% ACKNOWLEDGEMENTS -- Nothing to do here except comment out if you don't want it.
\begin{acknowledgements}
 	\subsection*{Personal}

This is where you thank your advisor, colleagues, and family and friends.

Lorem ipsum dolor sit amet, consectetur adipiscing elit. Vestibulum feugiat et est at accumsan. Praesent sed elit mattis, congue mi sed, porta ipsum. In non ullamcorper lacus. Quisque volutpat tempus ligula ac ultricies. Nam sed erat feugiat, elementum dolor sed, elementum neque. Aliquam eu iaculis est, a sollicitudin augue. Cras id lorem vel purus posuere tempor. Proin tincidunt, sapien non dictum aliquam, ex odio ornare mauris, ultrices viverra nisi magna in lacus. Fusce aliquet molestie massa, ut fringilla purus rutrum consectetur. Nam non nunc tincidunt, rutrum dui sit amet, ornare nunc. Donec cursus tortor vel odio molestie dignissim. Vivamus id mi erat. Duis porttitor diam tempor rutrum porttitor. Lorem ipsum dolor sit amet, consectetur adipiscing elit. Sed condimentum venenatis consectetur. Lorem ipsum dolor sit amet, consectetur adipiscing elit.

Aenean sit amet lectus nec tellus viverra ultrices vitae commodo nunc. Mauris at maximus arcu. Aliquam varius congue orci et ultrices. In non ipsum vel est scelerisque efficitur in at augue. Nullam rhoncus orci velit. Duis ultricies accumsan feugiat. Etiam consectetur ornare velit et eleifend.

Suspendisse sed enim lacinia, pharetra neque ac, ultricies urna. Phasellus sit amet cursus purus. Quisque non odio libero. Etiam iaculis odio a ex volutpat, eget pulvinar augue mollis. Mauris nibh lorem, mollis quis semper quis, consequat nec metus. Etiam dolor mi, cursus a ipsum aliquam, eleifend venenatis ipsum. Maecenas tempus, nibh eget scelerisque feugiat, leo nibh lobortis diam, id laoreet purus dolor eu mauris. Pellentesque habitant morbi tristique senectus et netus et malesuada fames ac turpis egestas. Nulla eget tortor eu arcu sagittis euismod fermentum id neque. In sit amet justo ligula. Donec rutrum ex a aliquet egestas.

\subsection*{Institutional}

If you want to separate out your thanks for funding and institutional support, I don't think there's any rule against it.  Of course, you could also just remove the subsections and do one big traditional acknowledgement section.

Lorem ipsum dolor sit amet, consectetur adipiscing elit. Ut luctus tempor ex at pretium. Sed varius, mauris at dapibus lobortis, elit purus tempor neque, facilisis sollicitudin felis nunc a urna. Morbi mattis ante non augue blandit pulvinar. Quisque nec euismod mauris. Nulla et tellus eu nibh auctor malesuada quis imperdiet quam. Sed eget tincidunt velit. Cras molestie sem ipsum, at faucibus quam mattis vel. Quisque vel placerat orci, id tempor urna. Vivamus mollis, neque in aliquam consequat, dui sem volutpat lorem, sit amet tempor ipsum felis eget ante. Integer lacinia nulla vitae felis vulputate, at tincidunt ligula maximus. Aenean venenatis dolor ante, euismod ultrices nibh mollis ac. Ut malesuada aliquam urna, ac interdum magna malesuada posuere.
\end{acknowledgements}

%%%%% ABSTRACT -- Nothing to do here except comment out if you don't want it.
\begin{abstract}
	Your abstract text goes here.  Check your departmental regulations, but generally this should be less than 300 words.  See the beginning of Chapter~\ref{ch:2-litreview} for more.

Lorem ipsum dolor sit amet, consectetur adipiscing elit. Pellentesque sit amet nibh volutpat, scelerisque nibh a, vehicula neque. Integer placerat nulla massa, et vestibulum velit dignissim id. Ut eget nisi elementum, consectetur nibh in, condimentum velit. Quisque sodales dui ut tempus mattis. Duis malesuada arcu at ligula egestas egestas. Phasellus interdum odio at sapien fringilla scelerisque. Mauris sagittis eleifend sapien, sit amet laoreet felis mollis quis. Pellentesque dui ante, finibus eget blandit sit amet, tincidunt eu neque. Vivamus rutrum dapibus ligula, ut imperdiet lectus tincidunt ac. Pellentesque ac lorem sed diam egestas lobortis.

Suspendisse leo purus, efficitur mattis urna a, maximus molestie nisl. Aenean porta semper tortor a vestibulum. Suspendisse viverra facilisis lorem, non pretium erat lacinia a. Vestibulum tempus, quam vitae placerat porta, magna risus euismod purus, in viverra lorem dui at metus. Sed ac sollicitudin nunc. In maximus ipsum nunc, placerat maximus tortor gravida varius. Suspendisse pretium, lorem at porttitor rhoncus, nulla urna condimentum tortor, sed suscipit nisi metus ac risus.

Aenean sit amet enim quis lorem tristique commodo vitae ut lorem. Duis vel tincidunt lacus. Sed massa velit, lacinia sed posuere vitae, malesuada vel ante. Praesent a rhoncus leo. Etiam sed rutrum enim. Pellentesque lobortis elementum augue, at suscipit justo malesuada at. Lorem ipsum dolor sit amet, consectetur adipiscing elit. Praesent rhoncus convallis ex. Etiam commodo nunc ex, non consequat diam consectetur ut. Pellentesque vitae est nec enim interdum dapibus. Donec dapibus purus ipsum, eget tincidunt ex gravida eget. Donec luctus nisi eu fringilla mollis. Donec eget lobortis diam.

Suspendisse finibus placerat dolor. Etiam ornare elementum ex ut vehicula. Donec accumsan mattis erat. Quisque cursus fringilla diam, eget placerat neque bibendum eu. Ut faucibus dui vitae dolor porta, at elementum ipsum semper. Sed ultrices dui non arcu pellentesque placerat. Etiam posuere malesuada turpis, nec malesuada tellus malesuada.
\end{abstract}

%%%%% MINI TABLES
% This lays the groundwork for per-chapter, mini tables of contents.  Comment the following line
% (and remove \minitoc from the chapter files) if you don't want this.  Un-comment either of the
% next two lines if you want a per-chapter list of figures or tables.
\dominitoc % include a mini table of contents
%\dominilof  % include a mini list of figures
%\dominilot  % include a mini list of tables

% This aligns the bottom of the text of each page.  It generally makes things look better.
\flushbottom

% This is where the whole-document ToC appears:
\tableofcontents

\listoffigures
	\mtcaddchapter
% \mtcaddchapter is needed when adding a non-chapter (but chapter-like) entity to avoid confusing minitoc

% Uncomment to generate a list of tables:
%\listoftables
%	\mtcaddchapter

%%%%% LIST OF ABBREVIATIONS
% This example includes a list of abbreviations.  Look at text/abbreviations.tex to see how that file is
% formatted.  The template can handle any kind of list though, so this might be a good place for a
% glossary, etc.
% First parameter can be changed eg to "Glossary" or something.
% Second parameter is the max length of bold terms.
\begin{mclistof}{List of Abbreviations}{3.2cm}

\item[1-D, 2-D] One- or two-dimensional, referring in this thesis to spatial dimensions in an image.

\item[Otter] One of the finest of water mammals.

\item[Hedgehog] Quite a nice prickly friend.

\end{mclistof} 


% The Roman pages, like the Roman Empire, must come to its inevitable close.
\end{romanpages}


%%%%% CHAPTERS
% Add or remove any chapters you'd like here, by file name (excluding '.tex'):
\flushbottom
% \begin{savequote}[8cm]
% \textlatin{Neque porro quisquam est qui dolorem ipsum quia dolor sit amet, consectetur, adipisci velit...}

% There is no one who loves pain itself, who seeks after it and wants to have it, simply because it is pain...
%   \qauthor{--- Cicero's \textit{de Finibus Bonorum et Malorum}}
% \end{savequote}

\chapter{Introduction}\label{ch:1-intro}

% \minitoc

\section{Motivation}
Quantum circuit compilation is one of the oldest problems in quantum computing.
The problem is to find a sequence of elementary gates that implements a given unitary transformation.
Often stated as a constrained optimisation problem: find a ``quantum circuit'' implementing
the desired computation that minimises a cost function such as the number of gates
of type X, or the \emph{depth} of the computation.

As quantum computers become bigger and computational demands from users increase,
new challenges are emerging in quantum compilation 
\begin{itemize}
  \item Circuits get bigger, making the compilation more computationally expensive.
  \item To scale and improve performance, quantum hardware architecture becomes more complex.
  The result is more special-purpose gate sets, additional execution constraints and more edge cases
  in the compilation problem.
  \item Quantum supercomputing: given the cost and scarcity of quantum computational
  resources, quantum compilation will increasingly be expected to leverage large
  classical compute clusters to lighten as much as possible the quantum resource requirements.
\end{itemize}
The goal of this thesis is to explore new research questions that arise from these developments.
All of these changing demands also come in hand in hand with substantial engineering challenges.

\section{Contributions}
The work we present here contributes towards the scaling up of quantum compilation
in three ways.
\begin{description}
  \item[Quantum Supercomputing.]
    We explore the new challenges for compilation that will arise from larger quantum
    computations, and in particular from the inevitable merging of quantum and classical
    large compute clusters. We explore current tools, their limitations and the
    requirements for future solutions. See \cref{ch:3-quantumsuper}.
  \item[Novel solutions to compilation.]
    Building on a proven \emph{supercompilation} approach, we 
    present two novel results that significantly improve the viability and performance 
    of rewrite rule-based quantum optimisation.
    The first is a fast pattern matching algorithm (\cref{ch:4-pattern}).
    This is semantics-agnostic, meaning that it can be applied not only on quantum circuits
    but also on hybrid quantum-classical computations and even pure-classical Machine Learning
    model training.
    The second is a contribution to circuit optimisation using the ZX-calculus, where we show
    how the pattern matching above can be used in causal-flow preserving ZX optimisations (\cref{ch:5-zx}). 
  \item[Tooling for large scale compilation.]
    We make performance-focused open source infrastructure available that can be
    used to scale compilation to large quantum computations.
    Significant enginnering efforts have gone into the design and implementation of
    this tooling, which we describe in \cref{ch:6-futcompiler}.
\end{description}
\begin{savequote}[8cm]
Alles Gescheite ist schon gedacht worden.\\
Man muss nur versuchen, es noch einmal zu denken.

All intelligent thoughts have already been thought;\\
what is necessary is only to try to think them again.
  \qauthor{--- Johann Wolfgang von Goethe \cite{von_goethe_wilhelm_1829}}
\end{savequote}

\chapter{\label{ch:2-litreview}Background}

% \minitoc

\chapter{Quantum Supercomputing}\label{ch:3-quantumsuper}
\chapter{Pattern Matching (see separate paper)}\label{ch:4-pattern}
\chapter{Causal ZX Optimisation (next paper tbd)}\label{ch:5-zx}
\chapter{Large Scale Quantum Compilation Platform}\label{ch:6-futcompiler}
\chapter{Conclusion}\label{ch:7-conclusion}


%% APPENDICES %% 
% Starts lettered appendices, adds a heading in table of contents, and adds a
%    page that just says "Appendices" to signal the end of your main text.
\startappendices
% Add or remove any appendices you'd like here:
% \include{text/appendix-1}


%%%%% REFERENCES

% JEM: Quote for the top of references (just like a chapter quote if you're using them).  Comment to skip.
% \begin{savequote}[8cm]
% The first kind of intellectual and artistic personality belongs to the hedgehogs, the second to the foxes \dots
%   \qauthor{--- Sir Isaiah Berlin \cite{berlin_hedgehog_2013}}
% \end{savequote}

\setlength{\baselineskip}{0pt} % JEM: Single-space References

{\renewcommand*\MakeUppercase[1]{#1}%
\printbibliography[heading=bibintoc,title={\bibtitle}]}


\end{document}
